% TOP_PROBLEM: {"problemId": "problem.arnold-givental-lagrangian-intersection-conjecture", "canonicalTitle": "Arnold-Givental Lagrangian intersection conjecture", "releaseRank": 395, "primaryDomain": "geometry_topology", "primaryDomainLabel": "Geometry & topology", "displayStatus": "open_with_solved_subcases", "releaseStatus": "open_with_solved_subcases"}
% !TeX program = lualatex
% ATTEMPT_STATUS: unresolved
\documentclass[11pt]{article}
\usepackage[margin=25mm]{geometry}
\usepackage{fontspec}
\setmainfont{Latin Modern Roman}
\usepackage{amsmath,amssymb,mathtools}
\usepackage{unicode-math}
\setmathfont{Latin Modern Math}
\usepackage{xurl,hyperref}
\hypersetup{colorlinks=true,urlcolor=blue}
\setcounter{secnumdepth}{0}
\setlength{\parindent}{0pt}
\setlength{\parskip}{5pt}
\setlength{\emergencystretch}{3em}
\begin{document}
\section*{\texorpdfstring{Arnold-Givental Lagrangian intersection conjecture}{Arnold-Givental Lagrangian intersection conjecture}}\label{395}
\subsection{Definitions and mathematical statement}
A real symplectic manifold is a compact symplectic $(M,\omega)$ with an involution $\tau$ satisfying $\tau^*\omega=-\omega$. Its nonempty fixed locus $L=\operatorname{Fix}(\tau)$ is Lagrangian: it has half the dimension of $M$ and $\omega|_{TL}=0$. A Hamiltonian diffeomorphism is the time-one map of a time-dependent vector field $X_{H_t}$ defined by $\iota_{X_{H_t}}\omega=-{\,\mathrm{d}} H_t$. If $L$ and $\phi(L)$ are transverse, the conjecture asserts
\[
 \#(L\cap\phi(L))\ge\sum_i\dim_{{\mathbb{F}}_2}H_i(L;{\mathbb{F}}_2).
\]
Compactness and transversality make the intersection finite. The real-locus hypothesis is part of the statement; it is not a lower bound for arbitrary Lagrangians or for non-Hamiltonian symplectomorphisms.
\par\smallskip\textit{Formulation and concept references:} \href{https://arxiv.org/html/1708.09628v1}{[S1]}.

\subsection{Short English statement}
Must every transverse Hamiltonian image of a compact real Lagrangian meet it in at least as many points as its total mod-two Betti number?
\subsection{Research attempt: exact graphs, the sphere, and sharp torus examples}
\textbf{Status.} The universal intersection bound is \textbf{UNRESOLVED} by this attempt. We prove it for sufficiently small Hamiltonian perturbations explicitly represented by exact one-form graphs, prove the equator case for every Hamiltonian diffeomorphism of the sphere, and construct transverse examples on symplectic tori attaining the full bound. The remaining problem is the passage from these models to unrestricted Hamiltonian isotopies without additional Floer-theoretic hypotheses.

\textbf{Bibliographic and scope audit.} The URL retained in [S1] actually leads to Frauenfelder's 2017 paper \emph{Smooth nondisplaceability for fixed point sets of involutions}, not to his 2004 paper \emph{The Arnold--Givental conjecture and moment Floer homology}. The retrieved 2017 paper states the selected bound as Conjecture 1.2 and proves a weaker nondisplaceability conclusion when the ambient Euler characteristic is odd. Its Theorem 1.1 is a smooth topological theorem. Neither it nor its Corollary 1.4 is the full Betti-number estimate. We preserve the catalogue entry and add the correct identification as [E1]. The arguments below do not turn nondisplaceability into a rank bound without an additional reason.

\subsubsection{1. Why the fixed locus is Lagrangian}
At a fixed point $x$, the derivative $A=d\tau_x$ satisfies $A^2=I$ and
$\omega(Au,Av)=-\omega(u,v)$. Its $+1$ and $-1$ eigenspaces $E_+$ and $E_-$ span $T_xM$. Each is isotropic: for $u,v$ in the same eigenspace the preceding identity reads $\omega(u,v)=-\omega(u,v)$.

The pairing between $E_+$ and $E_-$ is nondegenerate. If $u\in E_+$ pairs to zero with $E_-$, it already pairs to zero with $E_+$, hence with all of $T_xM$, so $u=0$. The same reasoning in the other direction shows $\dim E_+=\dim E_-=\tfrac12\dim M$. Using a $\tau$-invariant auxiliary Riemannian metric, the exponential chart at $x$ is equivariant for $\tau$; its fixed points near $x$ are exactly the image of $E_+$. Thus $L$ is a smooth half-dimensional submanifold and $\omega|_{TL}=0$.

This argument works component by component and does not assert that $L$ is connected or orientable. Mod-two homology is therefore particularly natural. Compactness makes $L$ a compact manifold with finitely many components, and its total mod-two Betti number
$B(L)=\sum_i\dim H_i(L;\mathbb F_2)$ is finite.

\subsubsection{2. An explicit local class in which Morse theory gives the bound}
Use the Lagrangian neighborhood theorem as an external structural input: a neighborhood of a compact Lagrangian $L$ is symplectomorphic to a neighborhood of the zero section in $T^*L$. In local coordinates $(q,p)$ take
$\omega=\sum_i dq_i\wedge dp_i$. For a smooth function $f:L\to\mathbb R$, the Hamiltonian $H(q,p)=\epsilon f(q)$ has, under the convention in the question,
\[
 \dot q=0,\qquad \dot p=\epsilon\,df(q).               \tag{1}
\]
Indeed contraction with $\omega$ gives $-\dot p\,dq=-\epsilon df$.

Choose a fiber cutoff which equals one in a neighborhood containing all segments $p=t\epsilon df(q)$, $0\le t\le1$, and has compact support in the chosen cotangent neighborhood. For sufficiently small $\epsilon$, the Hamiltonian extended by zero to $M$ is smooth and its flow from the zero section remains where the cutoff is identically one. Therefore its time-one image is exactly
\[
 L_f=\operatorname{graph}(\epsilon df).
\]
Its intersections with $L$ are precisely the critical points of $f$. In a chart at such a point, the tangent spaces are the horizontal space and the graph of $\epsilon\operatorname{Hess}f$. They are transverse exactly when the Hessian is nonsingular. Thus a Morse function yields a transverse Hamiltonian image and
\[
 \#(L\cap L_f)=\#\operatorname{Crit}(f)\ge B(L).        \tag{2}
\]
The last inequality follows from Morse theory with $\mathbb F_2$ coefficients: the Morse chain group has one generator for each critical point, and its homology is $H_*(L;\mathbb F_2)$. Algebraically, for any finite chain complex,
\[
 \dim C=\dim H(C)+2\dim\operatorname{im}\partial.        \tag{3}
\]
To verify (3), use $\dim C=\dim\ker\partial+\dim\operatorname{im}\partial$ and $\dim H=\dim\ker\partial-\dim\operatorname{im}\partial$, summing degrees. It applies also to disconnected and nonorientable $L$ over $\mathbb F_2$.

This proves a genuine family of Hamiltonian intersection inequalities, not just an infinitesimal count. However, the graph is exact because it was explicitly constructed from $f$. We have not asserted that every symplectic graph of a closed one-form is exact, or that an arbitrary global Hamiltonian image stays in one cotangent neighborhood along a suitable isotopy.

\subsubsection{3. The sphere equator: a global argument using area}
Let $M=S^2$ with its usual area form and let $\tau$ be reflection across the equatorial plane. This map reverses the area form, and its fixed locus is the equator $L\cong S^1$, with $B(L)=2$. Let $\phi$ be any Hamiltonian diffeomorphism, and put $C=\phi(L)$.

The curve $C$ is a smoothly embedded circle. It separates the sphere into the images of the northern and southern open hemispheres. Because Hamiltonian flow preserves the symplectic area, each component has area exactly half the total. If $C\cap L$ were empty, the connected curve $C$ would lie in one open hemisphere $U$. The Jordan curve theorem in this disk gives a component $D$ of $S^2\setminus C$ contained in $U$. Its area is strictly less than the area of $U$: the complement of the compact curve in $U$ has a nonempty open region on the side facing the equatorial boundary, and such a region has positive area. This contradicts the half-area conclusion for both components of $S^2\setminus C$. Hence $C$ meets $L$.

If the intersections are transverse, their number is even. To see this, travel once around $C$, marking whether the current point is in the northern or southern hemisphere. Each transverse crossing changes that sign, and at the end it must return to its initial sign. The nonzero even number is therefore at least two. This proves the exact conjectural bound for the equator, for every Hamiltonian diffeomorphism, without a smallness assumption.

Two limitations are visible in the proof. The equal-area property uses the particular real involution and the ambient symplectic area. A small embedded contractible circle not enclosing half the sphere need not be nondisplaceable. Also, the alternating-sign argument gives a numerical count because the complement has two distinguished components; there is no automatic higher-dimensional analogue giving $B(L)$.

\subsubsection{4. Sharp examples on a torus, including all real components}
On
\[
 M=\mathbb R^{2r}/\mathbb Z^{2r},\qquad
 \omega=\sum_{j=1}^r dq_j\wedge dp_j,
\]
let $\tau(q,p)=(q,-p)$. This is an anti-symplectic involution. Its fixed locus is
\[
 L=\coprod_{\eta\in\{0,1/2\}^r}\{p=\eta\},
\]
a union of $2^r$ copies of $\mathbb T^r$. The mod-two homology of $\mathbb T^r$ has total dimension $2^r$, by the product cell decomposition or the K\"unneth formula. Thus
\[
 B(L)=4^r.                                             \tag{4}
\]
It would be a factor-$2^r$ error to keep only the component $p=0$.

Take the globally defined Hamiltonian
\[
 H(q,p)=\epsilon\sum_{j=1}^r\cos(2\pi q_j),\qquad
 0<2\pi\epsilon<1/4.
\]
Equation (1) gives
\[
 \phi(q,p)=\bigl(q,p-2\pi\epsilon(\sin(2\pi q_j))_{j=1}^r\bigr).
                                                               \tag{5}
\]
Because each displacement has absolute value less than $1/4$, an image of one component can intersect a real component only when their $p$-labels are the same. Intersection then requires $\sin(2\pi q_j)=0$ for every $j$, so each $q_j$ is $0$ or $1/2$. Each component contributes exactly $2^r$ intersections and the total is $4^r$.

The Hessian is diagonal with entries
$-4\pi^2\epsilon\cos(2\pi q_j)$, each nonzero at those points. Hence all intersections are transverse. The conjectural lower bound is attained in every even ambient dimension by this family.

This does not prove the bound for an arbitrary Hamiltonian on the torus. It verifies the normalization of the right-hand side, the disconnected fixed-locus convention, the Hamiltonian sign, and the possibility of equality. The tensor-product count is exact and is checked for small $r$ in the saved finite diagnostics.

\subsubsection{5. Why symplectic displacement alone is not a counterexample}
Already on $\mathbb T^2$, the translation
$T_\delta(q,p)=(q,p+\delta)$ displaces the fixed locus $p\in\{0,1/2\}$ whenever $\delta\notin\{0,1/2\}$ modulo one. It preserves $dq\wedge dp$. To ensure that this is not inadvertently a Hamiltonian counterexample, we can prove directly that a nonintegral vertical translation is not Hamiltonian.

For a Hamiltonian isotopy $\phi_t$ with a periodic Hamiltonian $H_t(q,p)$, take its lift to $\mathbb R^2$ starting at the identity. The vertical displacement of a point is
\[
 \Delta p(z)=\int_0^1\partial_qH_t(\phi_t(z))\,dt.
\]
Average over one fundamental square. Since $\phi_t$ preserves area on the torus,
\[
 \int_{\mathbb T^2}\Delta p(z)\,dz
 =\int_0^1\int_{\mathbb T^2}\partial_qH_t(w)\,dw\,dt=0. \tag{6}
\]
If the endpoint were $T_\delta$, this continuous lift would differ from the standard translation lift by a single integer deck translation; its vertical displacement would be the constant $\delta+k$ for some integer $k$. Equation (6) would force $\delta+k=0$. For $\delta\notin\mathbb Z$ that is impossible.

Thus the visibly displacing translations lie outside the Hamiltonian class. The proof uses the endpoint, not merely the fact that the most obvious generating vector field has a nonexact contraction. A non-Hamiltonian choice of isotopy by itself would not exclude a different Hamiltonian isotopy with the same endpoint.

\subsubsection{6. The global Floer template and what it still needs}
The local Morse argument suggests replacing critical points by intersection points of $L$ and $\phi(L)$. If one had a finite chain complex over $\mathbb F_2$ with one generator at each transverse intersection and homology isomorphic to $H_*(L;\mathbb F_2)$, then (3) would give the conjecture immediately. This is a sufficient algebraic template, not a construction supplied in this attempt.

The analytic construction normally counts suitable pseudoholomorphic strips with boundaries on the two Lagrangians. To make such a count a differential and prove invariance, one must control compactifications of one-dimensional moduli spaces, including disk and sphere bubbling, and identify the resulting homology. The real involution provides symmetries, but saying that disks occur in reflected pairs does not by itself prove cancellation: one must check fixed configurations, regularity, boundary strata and compatibility with the coefficient system.

A purely algebraic stress test shows why a weaker invariant is insufficient. Euler characteristic alone determines
$\sum_i(-1)^i\dim H_i$, whereas the desired bound uses the unsigned sum. For $L=S^1$, Euler characteristic is zero while $B(L)=2$; for $\mathbb T^r$ with $r>0$, it is also zero while $B(L)=2^r$. Even a nonzero self-intersection or an odd-parity obstruction generally gives only at least one intersection. It does not supply the chain homology needed for (3).

Similarly, a statement that every image intersects $L$ cannot be amplified to $B(L)$ by perturbing the image repeatedly. The intersections for different perturbations need not be distinct points of the original image. The source's nondisplaceability theorem is therefore useful context but not a replacement for the missing rank calculation.

\subsubsection{7. Adversarial verification and the remaining question}
All local graph intersections use an honest compactly supported Hamiltonian after cutoff, and the flow stays in its constant-cutoff region. Compactness guarantees such a common region for the bounded set of covectors $\epsilon df$. Morse transversality is checked by the Hessian rather than assumed from a count. The sphere argument uses both area equality and transversality; the count without transversality would need a different formulation. The torus family includes every connected component and excludes cross-component intersections by a strict displacement bound.

The finite diagnostic checks the component labels, the $2^r$ critical points per component, the nonzero Hessian signs and the equality $\#(L\cap\phi(L))=4^r$ for the explicit models. It cannot validate the analytic compactness or continuation maps for an arbitrary real Lagrangian.

The strongest results obtained are the exact-graph subcase, the global equator theorem and the sharp torus construction. The first unresolved general step is a well-defined intersection complex with the required homology, or another argument giving its full rank consequence, for every real fixed locus in the catalogue without additional positivity or bubbling restrictions. The selected universal inequality remains \textbf{UNRESOLVED} by this attempt.

\subsection{Sources}
\begin{itemize}
\item[S1]Urs Frauenfelder, Arnold-Givental conjecture and moment Floer homology, 2017.
\url{https://arxiv.org/html/1708.09628v1}.
\textit{Formulation source.}
\item[E1]U. Frauenfelder, \emph{Smooth nondisplaceability for fixed point sets of involutions}, arXiv:1708.09628v1 (2017), Conjecture 1.2, Theorem 1.1 and Corollary 1.4.
\url{https://arxiv.org/html/1708.09628v1}.
\textit{Correct identification of the URL printed in S1; full text inspected 27 September 2026. The distinct moment Floer paper appears as reference 4 there and dates to 2004.}
\end{itemize}
\end{document}
